\documentclass[12pt,a4paper]{article}

% ------------------------------------------------
% PAQUETES
% ------------------------------------------------
\usepackage[spanish, es-nodecimaldot]{babel}
\usepackage[utf8]{inputenc}
\usepackage[T1]{fontenc}

\usepackage[
    top=2.5cm,
    bottom=2.5cm,
    left=3cm,
    right=3cm
]{geometry}

\usepackage{setspace}
\usepackage{parskip}
\usepackage{amsmath}
\usepackage{lmodern}

% ------------------------------------------------
% FORMATO DEL DOCUMENTO
% ------------------------------------------------
\onehalfspacing

\setlength{\parindent}{0pt}
\setlength{\parskip}{0.5em}

% ------------------------------------------------
% INFORMACIÓN
% ------------------------------------------------
\title{\textbf{Huella isotópica del vapor atmosférico en Salamanca, Guanajuato:\\
develando el impacto de la contaminación antropogénica}}

\author{Brenda Avelina López Muñiz}
\date{\today}

% ------------------------------------------------
% DOCUMENTO
% ------------------------------------------------
\begin{document}

\maketitle

\begin{center}
\textit{Maestría en Ciencias la Tierra, UNAM}
\end{center}

\vspace{0.5cm}

El agua atmosférica constituye un componente fundamental del ciclo 
hidrológico y su composición isotópica puede proporcionar información
sobre los procesos físicos y químicos que intervienen en su formación,
transporte y transformación. Los isótopos estables del hidrógeno y 
del oxígeno, expresados mediante $\delta$D y $\delta^{18}$O, han
sido ampliamente utilizados como trazadores de procesos hidrológicos
debido a que su composición puede modificarse durante fenómenos 
como la evaporación, condensación, precipitación y mezcla de 
diferentes fuentes de humedad. Sin embargo, la composición 
isotópica del vapor atmosférico también puede estar influenciada 
por procesos asociados a actividades antropogénicas, entre ellos 
la combustión de combustibles fósiles.

La presente investigación busca caracterizar la huella isotópica del 
vapor atmosférico en Salamanca, Guanajuato, con el propósito de explorar
la posible influencia de fuentes antropogénicas sobre la composición isotópica 
del agua atmosférica. Salamanca constituye un sitio de interés debido a la
presencia de diversas actividades industriales y urbanas que contribuyen a la 
emisión de contaminantes atmosféricos. En este contexto, el estudio integra
información isotópica con datos de calidad del aire y variables meteorológicas
para analizar las posibles relaciones entre la composición del vapor atmosférico
y las condiciones ambientales presentes durante el periodo de muestreo.

Para el desarrollo de la investigación se realizaron muestreos de agua atmosférica 
mediante un sistema de captación y condensación, obteniéndose muestras representativas
del vapor presente en el ambiente. La composición isotópica de las muestras se determinó 
mediante $\delta$D y $\delta^{18}$O y posteriormente se comparó con respecto a la 
línea meteórica regional, mundial y la huella isotópica del vapor de 
combustión. De manera complementaria, se desarrolló un 
experimento bajo condiciones ambientales controladas, utilizando agua destilada como 
fuente de humedad a una temperatura de $30\,^{\circ}\mathrm{C}$ y una humedad 
relativa de 50\,\%. Este experimento permitió estudiar la modificación de la 
composición isotópica durante los procesos de evaporación y condensación y 
establecer una referencia para interpretar el comportamiento de las muestras 
ambientales.

El análisis se complementa con información de contaminantes atmosféricos, 
incluyendo CO, NO, NO$_x$, NO$_2$, O$_3$, SO$_2$, PM$_{10}$ y PM$_{2.5}$, 
así como variables meteorológicas como velocidad y dirección del viento, 
temperatura, humedad relativa, presión atmosférica y radiación ultravioleta. 
Estos datos permiten contextualizar temporal y espacialmente las muestras 
isotópicas y explorar la relación entre las condiciones de contaminación, la 
dinámica atmosférica y la composición del vapor de agua. Asimismo, el empleo de
trayectorias de masas de aire a través de un modelo de transporte
y dispersión permite considerar el origen y evolución de las
masas de aire que alcanzan la zona de estudio.

La integración de estas aproximaciones busca determinar si las variaciones 
observadas en la composición isotópica del vapor atmosférico pueden explicarse 
mediante procesos naturales del ciclo hidrológico o si presentan características
compatibles con la influencia de fuentes antropogénicas. De esta manera, la 
investigación pretende contribuir al uso de los isótopos estables como una 
herramienta complementaria para el estudio de la contaminación atmosférica y 
de las interacciones entre las actividades humanas y el ciclo hidrológico. Los 
resultados podrían servir como base para desarrollar, en etapas posteriores, 
indicadores isotópicos que permitan evaluar la influencia de fuentes de emisión
sobre el vapor atmosférico en zonas urbanas e industriales.

\end{document}
```
